\chapter{Scope}
\label{chap:scope}

% Closed question: What kind of claim can a name carry?

A class name is where the argument becomes sayable.

Before any airliner leaves the gate, the crew runs through a list. The
list has names: flaps, trim, fuel pumps, transponder, clearance, runway,
weather, final readiness. Each name is brief. Each name corresponds to a
state the aircraft must be in. The crew calls the name; the responder
confirms or withholds. The takeoff proceeds only after the list closes.
If a name is called and the corresponding state cannot be confirmed, the
aircraft does not take off. The list is short. The aircraft is enormous.
The list controls the aircraft.

A passenger who has never flown will still recognize the scene. The
checklist is famous. Its discipline is famous. The reader does not have
to be told that the checklist matters; the reader knows that two hundred
people are about to be lifted into the sky. The checklist is the
smallest visible apparatus inside a much larger apparatus the reader
already trusts.

The aviation industry's adoption of the checklist is not an old
tradition. It is the result of a specific accident, the Boeing Model 299
crash at Wright Field in 1935, in which a pilot took off without
releasing the elevator lock. Two crew members died. The investigators
concluded that the airplane had become too complicated for any single
pilot to fly from memory. The fix was not to demand smarter pilots. The
fix was a list of named items the crew would call out before takeoff.
The names of the items were short. The names were specific. The names
were called in a fixed order. The list closed when every item had been
confirmed.

The checklist did not make the airplane easier to fly. It made the
airplane safer to fly by externalizing what previously sat in one
person's head. The names made the procedure inspectable. A copilot
could check the captain. An observer could check both. The fix was not
technical. It was rhetorical. The names did the work that no individual
memory could be trusted to do.

The project's class names share that structure. They externalize what
the formal record must do before the next move is admissible. They make
the procedure inspectable. A reader can check a writer. A second reader
can check both. The names do work that no individual claim could be
trusted to do on its own.

The procession of class names is longer than the takeoff checklist. The
names are stranger. The discipline is the same.

\section{The Name As Instrument}
\label{se:the-name-as-instrument}

People trust names every day before they trust proofs.

A runway number, a calibration sticker, a boarding pass, a courtroom
role, a warning placard, a medication label, a fuse rating: each name
permits or forbids action. The name is not the whole truth about the
object. It does not explain the object's history or contents in full.
It changes what can be done next. A boarding pass with the wrong gate
number sends the passenger to the wrong airplane. A fuse with the wrong
rating burns the circuit. The name controls what happens after.

The use of names as instruments is older than industrial design. A
street name does not contain the street; it does not describe the
buildings, the residents, or the traffic patterns. A street name
controls what the postal service does next. A package addressed to a
street name will be delivered to the buildings along that street. A
package addressed to an unnamed location will return to the sender. The
name is the gate the postal system passes through. The postal system
has elaborate internal apparatus, but the name is the public surface
the apparatus reads.

A medical wristband on a patient in a hospital is another instance. The
band carries a name, a barcode, and a number. The name is short enough
to read aloud. The barcode is precise enough to scan. The number is
unique enough to identify. Each is an instrument layered over the same
underlying patient. The nurse who administers medication does not
consult the patient's complete medical record at each dose. The nurse
reads the wristband, confirms the name, scans the barcode, and verifies
that the medication's order matches the patient's record. The wristband
is the gate. The full medical record is what the wristband points at.

A class name in formal code works the same way. The name does not
contain the whole specification of the obligation. It points to the
obligation. Code that invokes the class name promises to respect what
the class requires. Code that omits the class name has not made the
promise. The compiler then refuses the next move. The name is the
smallest public surface of the obligation. The obligation itself lives
in the class definition. The class definition is long. The class name
is short. Shortness is what makes the name usable as an instrument.

In the cockpit, the captain calls flaps. The first officer reads the
flap-position indicator. The flap-position indicator is itself an
instrument, with internal mechanism, calibration history, and certified
accuracy. The captain does not call the mechanism. The captain calls
flaps. The single word is enough because the responder knows what flaps
means inside this cockpit on this aircraft on this morning. The single
word is the gate through which the mechanism's reading enters the
takeoff record. The actual flap angle has been verified against a
procedure the airline has codified, a procedure the aircraft
manufacturer has specified, and a procedure the regulatory authority
has certified. The reading is layered. The name calls all the layers at
once.

A name as instrument is short. A long name cannot be called under
pressure. The pre-flight checklist runs in two minutes or less; a
checklist that took fifteen minutes would be skipped. A class name that
took a sentence to write would be substituted with a shorter alias.
Shortness is structural. The instrument has to fit in the moment in
which the instrument is needed.

A name as instrument is unambiguous within its context. A name that
could mean two states cannot be confirmed by a single response. The
captain calls flaps and the first officer answers set. There is no
question about which state set refers to. The cockpit has agreed in
advance that flaps means a specific control surface, and set means a
specific position of that surface, and the agreement is unambiguous to
the two people in the cockpit. \texttt{DISTINGUISHABLE} does not float
between several possible obligations. It names one specific obligation,
inside one specific formal vocabulary, and the obligation is the same
every time the class name is invoked.

A name as instrument is binding. Once the name is called and confirmed,
the takeoff moves on; the state behind the name is now treated as
established. The flap-position indicator does not get a second
verification at runway-entry. The first officer's set is taken as
licensing the next checklist item. If the flap position is wrong, the
failure surfaces somewhere else in the system: the takeoff performance
calculation, the stall warning, the post-incident analysis. Within the
checklist, set is binding. Class names are similarly binding. Once an
instance of \texttt{DISTINGUISHABLE} has been provided, subsequent code
may assume the witness is in place. The compiler does not re-check the
witness at each use. The provision is binding for the scope of the
use.

An instrument can be called and lie. A flap-position indicator can be
miscalibrated. A medical wristband can be on the wrong patient. A class
instance can be provided that satisfies the name's interface without
satisfying the obligation the name was meant to carry. Instrument-
failure is a real problem and a separate question from what the name
does when it works. The procession in Chapter 2 will return to
instrument-failure in several specific places.

A name's instrument function depends on a trained community. The
captain and first officer are trained to use the checklist. Pilots in
training rehearse the checklist hundreds of times before they are
licensed. Class names depend on a different trained community:
programmers who know what \texttt{DISTINGUISHABLE} demands, readers
who have learned the project's vocabulary, and reviewers who can spot
a class name invoked without the obligation backing it.

A name can also fail as an instrument by being overused. A boy-who-cried-
wolf name no longer controls behavior. A warning placard posted
everywhere is read nowhere. A class name applied to every record loses
its discriminating power. The project's class names are used precisely.
\texttt{DISTINGUISHABLE} is not invoked on records that do not need the
obligation. \texttt{REPEATABLE} is not invoked on processes that do
not need the obligation. The instrument function depends on the
obligation being a real demand the name filters for. A name that
filters for nothing is not an instrument; it is decoration.

Instrument is the visible surface. The next question is what the
surface controls access to.

\section{The Name As Obligation}
\label{se:the-name-as-obligation}

A good name makes evasion harder.

A record that wears \texttt{REPEATABLE} owes a second reading. The
reading must happen again, under the same conditions, with the same
result. A single reading does not satisfy the name. A reading made
under different conditions does not satisfy the name. The name is
short, but its short shape forces a specific commitment. A process
that has been observed once is not \texttt{REPEATABLE} yet, no matter
how confidently the observer recalls it. The name will not let the
observation through until the second observation, under the same
conditions, has been made.

A record that wears \texttt{HALTED} owes a stopped process. The process
must have stopped at a recognizable output. A process that might still
produce more information does not satisfy the name. A process that has
been declared finished without an output cannot wear the name. The
short word refuses a class of weak claims. A scientific study cannot
be \texttt{HALTED} while data collection continues. A computation
cannot be \texttt{HALTED} while it is still running. The name is short
and the refusal it carries is sharp.

A record that wears \texttt{INFERRED} owes a record-permitted
conclusion. The conclusion must be the kind a disciplined record
permits. A conclusion reached by status, intuition, or institutional
pressure does not earn the name. The name does not soften the claim;
it narrows it. A sentence that wears \texttt{INFERRED} must point to
the record that licenses it. A sentence that wears \texttt{INFERRED}
without pointing to the record is asking the name to do work it has
refused to do.

Class names are not labels glued to internal structure. They are
obligations the structure has had to take on. A record carrying the
name \texttt{REPEATABLE} has accepted the obligation to be
reproducible. A record carrying the name \texttt{ADMISSIBLE} has
accepted the obligation to have entered through a gate. A record
carrying the name \texttt{WITNESSED} has accepted the obligation to
have been transported through a recordable path. The externalization
makes the obligations checkable.

In the cockpit, the captain calls fuel checked. The first officer reads
the fuel quantity, the planned reserve, and the diversion margin. Each
reading is an obligation. The crew does not get to nod through fuel
checked. The phrase requires a quantity in pounds or kilograms or
liters, a unit, a planned consumption for the flight, a reserve for
unexpected diversions, and a comparison against minimum legal reserves
for the route. If the quantity is below the planned departure
threshold, the call cannot be confirmed. The takeoff stops. The
obligation is inspectable because the name is short enough to be heard
and specific enough to be checked. The obligation has been codified by
the airline, by the aircraft manufacturer, and by the regulatory
authority. The captain and first officer did not invent the obligation;
they inherited it and they answer to it.

The codification makes the obligation enforceable. A name that points
at an obligation no one has written down is not enforceable. The
captain who decides personally what fuel checked should mean has not
made the call. The call is what the airline, the manufacturer, and the
regulator together have decided fuel checked must mean. The captain's
role is to make the call honestly under that codification. The
codification preceded the captain. The codification will outlive this
morning's flight.

The obligation each class names is fixed by the definition, not
invented in the moment. A reader can look up what
\texttt{DISTINGUISHABLE} requires, by consulting the class definition
in the formal volumes. The obligation is publicly held: any code that
claims an instance of the class can be inspected against the
definition. A reviewer can ask whether the instance actually satisfies
the obligation, or whether the instance is a placeholder pretending to
satisfy it. The obligation is small enough to fail: the compiler can
refuse the instance. The class name is the public form of an
obligation that the code has had to accept in order to proceed.

A class name carries one obligation. The next class name in the
procession carries the next obligation. The procession's force comes
from the fact that the obligations build on one another.
\texttt{ADMISSIBLE} does not satisfy \texttt{REPEATABLE}.
\texttt{REPEATABLE} does not satisfy \texttt{WITNESSED}. Each
obligation must be satisfied in order, and the satisfaction of one
does not waive the others. The class names are independent
obligations stacked in a forced sequence.

A name as obligation has three properties to hold. The obligation is
named in advance, by a codification that precedes the moment of use.
The obligation is publicly held, so that a reviewer can check the
claim against the codification. The obligation is small enough to
fail, so that the codification can refuse a candidate instance that
does not meet the demand. A purported obligation that lacks any of
these is not yet an obligation. It is an aspiration.

The class names of the project accept all three properties. The
checklist callouts of the pre-flight procedure accept all three
properties. The medical wristband's identification accepts all three
properties. The street name on the postal label accepts all three
properties. Named obligation is everywhere in ordinary public life.
The project's class names sit inside that grammar.

The individual obligations are independently inspectable. The next
question is what gives the sequence rhetorical weight.

\section{The Name As Public Argument}
\label{se:the-name-as-public-argument}

A private proof can be correct and still fail to persuade the public.

A mathematician may have an airtight proof that a function converges.
The mayor of a small town cannot read the proof. The mayor must still
decide whether to license the bridge that the function describes. The
mathematician and the mayor need a shared interface. The interface is
not the proof; it is the public claim the proof licenses. The
mathematician says, after this proof, the bridge load is admissible at
the stated rating. The mayor can use that sentence.

The shared interface is not a simplification of the proof. The proof
remains intact. The proof's authority does not transfer to the public
sentence by accident. It transfers because the public sentence carries,
in compressed form, the obligations that the proof's discipline
imposed. The sentence is short because the obligations have already
been satisfied by the proof. The sentence is honest because the
sentence does not promise more than the proof's obligations covered.
The mayor who licenses the bridge under that sentence is not licensing
the bridge on the strength of the mathematician's reputation. The
mayor is licensing the bridge on the strength of a named obligation:
load rating, tested under specified conditions, admissible at the
stated value.

The mathematician's proof remains intact behind the sentence. What the
public hears is a short sentence whose obligations the public can
recognize. Load rating. Tested sample. Engineering standard. Each is a
name with an inspectable meaning. Together they form a small public
argument that the mayor, the city council, and the townspeople can use
as the basis for a decision.

The mayor's decision is not the mathematician's proof. The decision is
the public action the proof licenses, under the public argument the
proof produces. The two registers are connected by the named
obligations. The argument moves between them only because the
obligations have been preserved across the move. A bridge-rating
sentence that did not point at any obligation would not move the
mayor's decision; it would be a slogan. A bridge-rating sentence that
pointed at the load test, the engineering standard, and the
manufacturer's certification is doing public-argument work because the
obligations are named and inspectable.

The passenger does not hear the whole flight-control system before
takeoff. The passenger does not hear how flaps are
electromechanically actuated, how the transponder negotiates with
air-traffic control, how the weather briefing has been derived from
upper-atmosphere sounding balloons. The passenger hears that the
captain has run the checklist. The captain says cleared for takeoff
after the gate has closed. The passenger uses that sentence. The
sentence is short. It is also a public argument: the checklist has
closed, the gates have held, the aircraft may move.

The argument is not made by the captain alone. The captain is the
visible speaker. The argument is made by the checklist, by the
codification behind each item, by the training that lets the captain
and first officer execute the checklist accurately, by the maintenance
organization that prepared the aircraft, by the air traffic control
system that issued the clearance, and by the regulatory authority
that certified all of the above. The captain is the last speaker in a
long chain of speakers. The cleared for takeoff sentence is the
chain's compressed conclusion. The passenger trusts the sentence
because the chain is real, even though the passenger cannot personally
inspect each link.

A sequence of names becomes a public argument when three conditions
hold. The names appear in a forced order. The captain does not call
cleared for takeoff before flaps. The order is not arbitrary; the
obligations build on one another. Cleared for takeoff depends on
having confirmed weather, runway, fuel, flaps, trim, and clearance,
in that order, because each later item presupposes the earlier items.
A clearance issued before fuel is checked is not operationally
meaningful. The order encodes the dependency structure.

Each name in the sequence is independently inspectable. A passenger
who wanted to could ask which names the captain called and what the
responder said. The information is recorded in the cockpit voice
recorder, the flight data recorder, and the maintenance log. Each
name's confirmation is a piece of recorded behavior that any
investigator can later examine. Each class name's instance is
recorded in the project's source code. Each instance can be
inspected against the class definition. The procession is not a
sealed performance; it is a sequence of inspectable acts.

The closing sentence summarizes the sequence without erasing it.
Cleared for takeoff is short because the checklist is long. The
short sentence is honest only because the long checklist held. A
captain who said cleared for takeoff without running the checklist
would be making the same words mean something different: an
unsupported claim rather than a public argument. The shortness of
the sentence is licensed by the length of the procession behind it.
Remove the procession and the sentence becomes mere assertion.

A procession can be wrong. The order encoded in the class names is
itself a choice that can be debated. A reader who thinks the
procession is wrong can argue against the order, propose a different
order, and engage with the reasoning that the current order encodes.
The project's procession is not infallible. It is one specific
ordering of obligations, chosen because each link's dependency on the
previous link could be made explicit. A different procession might
be equally valid for a different decision. The procession's defense
is not that it is the only one, but that it is internally coherent
and externally inspectable.

A public can misunderstand a name. A reader may hear \texttt{ADMISSIBLE}
as carrying a softer obligation than it actually demands. The
correction is teaching, not invention. The chapter that follows walks
the procession in detail, with each name receiving its own
development. Misunderstanding is a teaching problem.

A captain can falsify the checklist. A class-name claim can be made
without the obligation being satisfied. The defense is inspection. The
flight recorder catches falsified checklists. The reviewer catches
falsified class instances. The defense depends on the institutional
infrastructure that makes inspection possible. The class names are
not self-enforcing; they are inspectable only if someone takes the
trouble to inspect them. The procession is honest about what
honoring it requires, and what failing it would look like.

\begin{handoff}

The class names of the formal project do the same kinds of work as
the pre-flight checklist, in the same kind of order, but for a
different decision. The decision is whether a used car should be
bought at the offered price. The decision is local, financial, and
personal. The discipline is the same.

Each section in Chapter 2 opens with a public pressure, names the
class that answers it, explains why the name had to appear, shows the
running car-inspection example, shows an independent example from
another domain, and hands off to the next pressure. The procession is
forced. Each name makes the next one necessary.

The class-name procession is longer than the pre-flight checklist and
stranger. The names include some that are deliberately impolite. The
discipline is recognizable. A reader who heard flaps as a name doing
public work in the cockpit will hear \texttt{ADMISSIBLE} and
\texttt{WITNESSED} and \texttt{INFERRED} as names doing public work
in the same grammar.

The decision at the end of Chapter 2 is small. The buyer may infer
that the car is worth buying, may infer that the car is worth buying
only after a repair credit, or may infer that the buyer should walk
away. The smallness of the decision lets the procession do its work
without importing a contested public domain. A reader who can follow
the procession in this small case can carry it to larger cases.

\end{handoff}
